% 2026-11-23 W11 第 2 部。Michelman 404--410 依 PDF 頁圖核過（406、407 兩頁以圖核註 20--32）；註 12--19、33--36 依文字層。
% Michelman 註 12--19、33--36 所引 Post Constitutional Domains、Lefort、Rawls、Michelman Ethics 等原件本機未取得，一律轉引，不建 bib。
% 註 29 標明 Stanley v. Georgia 為 Marshall 大法官之意見，註 26 標明 Cruzan 為 Brennan 之不同意見，其餘註只列判決。
% Michelman 註 19 作 Rawls A Theory of Justice 85--86，轉引自 Michelman 註 19。Post 438 頁註 33 依 JSTOR 本文字層與頁圖。
\chapter{Post 之回應式民主及其程序性}
\label{ch:20_responsive}

\begin{qaquestion}
Michelman\ix{persons}{michelman@Michelman, Frank I.} 在 404--410 頁處理 Post\ix{persons}{post@Post, Robert} 的民主\ix{terms}{01@\zhlidxcat{民主理論}!民主}理論。
\begin{enumerate}
\tightlist
\item 回應式民主\ix{terms}{01@\zhlidxcat{民主理論}!回應式民主}的內容是什麼，Post 如何以 Lefort\ix{persons}{lefort@Lefort, Claude} 的話說明它全然程序性？
\item 純粹程序標準\ix{terms}{02@\zhlidxcat{正當性與合理歧見}!純粹程序標準}與程序獨立之正確標準\ix{terms}{02@\zhlidxcat{正當性與合理歧見}!程序獨立之正確標準}如何區分？
\item Brennan\ix{persons}{brennan@Brennan, William J., Jr.} 以人性尊嚴為核心的判決，為何無法由純粹程序的民主觀說明？
\item Michelman 如何論證 Post 的理論含有實質前提，仇恨言論之例、最高之法\ix{terms}{03@\zhlidxcat{憲政主義與憲法之效力基礎}!最高之法}與無窮後退\ix{terms}{03@\zhlidxcat{憲政主義與憲法之效力基礎}!無窮後退}在論證中各居何位？
\end{enumerate}
\end{qaquestion}

\section{回應式民主之內容}
\label{sec:20_responsive-1}

Michelman 在 404 頁以「responsive democracy」〔回應式民主\ix{terms}{01@\zhlidxcat{民主理論}!回應式民主}〕稱呼 Post 的理想：一組社會與法律條件，使每個願意參與的人，得以透過自己對「national identity」〔國族認同\ix{terms}{01@\zhlidxcat{民主理論}!國族認同}〕或「social order」〔社會秩序\ix{terms}{01@\zhlidxcat{民主理論}!社會秩序}〕之集體決定所作的貢獻，實際行使自治\ix{terms}{01@\zhlidxcat{民主理論}!自治}。\footcite[404 \& nn.12--13]{michelman1998}依註 13，這兩個詞出自 Post 的 Constitutional Domains 275、185 頁。Michelman 把它們理解為涵蓋一國基本的立法制度與實踐，並「further」〔進一步〕把這些制度與實踐所反映的主要政治文化傾向也算在內，例如美國人對個人堅毅與獨立、競爭、獎勵功績的偏好。這項擴張是 Michelman 自己的詮釋，他以「I further take them to include」標明。\footcite[404]{michelman1998}

在 Michelman 的轉述裡，Post 主張，國族認同與社會秩序若是在極其自由的公共論述\ix{terms}{01@\zhlidxcat{民主理論}!公共論述}中明顯地鍛造出來，居民便能以「owners」〔所有人〕的身分認同它，並對由它而生的一般社會規範負責，因為每個人都能在其中看到個人自由投入的偏好與意見之總和的真實而公平的呈現（註 14 指 Constitutional Domains 273、311--312 頁）。依此，民主的政治正當性\ix{terms}{02@\zhlidxcat{正當性與合理歧見}!政治正當性}在於「each individual's warranted sense of autonomous participation in the process of creating the social order in which we live」〔每個人對於自己自主參與創造我們所處社會秩序之過程，有其有依據的感受〕。\footcite[404 n.14]{michelman1998}

這項主張的出發點是自由主義者面對的困境：各人經驗與評價的深刻分歧使任何重大社會安排問題都無法達成共識，所以由每個人完全自我決定的政治無法達成。Michelman 轉述的回應式民主理論，以「unprejudiced participatory access for all to a ceaseless communicative process」〔所有人無偏見地參與一個不停運作的溝通過程〕作為次佳的替代：個人把自己的意見與偏好向量投入過程，由過程加總，暫時決定社會生活的基本安排。理論的至高承諾，是讓每個人對於自己有自主而有效的貢獻保有有依據的感受。Post 因此主張須嚴格維持「a structure of [political] communication」〔政治溝通的結構〕絕對不受限制，對任何人關於社會秩序所持的任何觀點開放。\footcite[404 \& n.15]{michelman1998}Post 的原話是：「[i]f the state were to forbid the expression of a particular idea, the government would become, with respect to individuals holding that idea, heteronomous and nondemocratic」〔國家若禁止表達某項觀念，對持該觀念的人而言，政府便成為他律而不民主的〕（註 16 指 Constitutional Domains 304 頁）。維持一個相當寬廣的、絕對不受限制之公共論述\ix{terms}{01@\zhlidxcat{民主理論}!不受限制之公共論述}領域，則在其中不斷鍛造的國族認同，可說是所有人自由選擇的；由這項國族認同產生並「reflect」〔反映〕它的社會法律規範，才可說與自我決定相容（註 17 指 274 頁）。\footcite[405 \& nn.16--17]{michelman1998}

\section{全然程序性之主張}
\label{sec:20_responsive-2}

Michelman 說，Post 把這套理論表述為「bottomlessly procedural」〔程序到底〕，不以任何實質基礎為依據，也不蘊含任何實質預設。Post 引 Claude Lefort\ix{persons}{lefort@Lefort, Claude} 的話為此背書，民主是「a regime founded upon the legitimacy of a debate as to what is legitimate and what is illegitimate---a debate which is necessarily without any guarantor and without any end」〔奠基於一場辯論之正當性\ix{terms}{02@\zhlidxcat{正當性與合理歧見}!正當性}的政體，辯論所論為何者正當、何者不正當，辯論必然沒有擔保者，也沒有終點〕。\footcite[405 n.18]{michelman1998}Michelman 此處所引是 Post 在 Constitutional Domains 186 頁所引的 Lefort，Democracy and Political Theory 39 頁（David Macey 英譯，1988）。Michelman 讀出的含義是，若人民真正自治，則程序性一路貫徹到底：辯論本身、辯論的框架與程序規範，皆只能是人民以民主方式造成的樣子。\footcite[405]{michelman1998}

\section{純粹程序標準與程序獨立之正確標準}
\label{sec:20_responsive-3}

Michelman 為說明 Post 等人為何採此看法，區分根本法\ix{terms}{03@\zhlidxcat{憲政主義與憲法之效力基礎}!根本法}之「rightness」〔正確性〕的兩種標準。\footcite[405]{michelman1998}

\begin{longtable}{@{}L{0.20\linewidth}L{0.40\linewidth}L{0.32\linewidth}@{}}
\caption{根本法正確性之兩種標準（Michelman 405 頁及註 19）。}\\
\toprule
標準 & Michelman 所舉例 & 與 Rawls\ix{persons}{rawls@Rawls, John} 分類之關係 \\
\midrule
\endfirsthead
\toprule
標準 & Michelman 所舉例 & 與 Rawls 分類之關係 \\
\midrule
\endhead
純粹程序標準\ix{terms}{02@\zhlidxcat{正當性與合理歧見}!純粹程序標準}（pure-procedural） & 根本法在準確表現每個人經適當表達的偏好或意見之總和時為正確 & 直接對應 Rawls 所稱 pure procedural justice，正義或正確性等同於履行程序並遵守其結果 \\
程序獨立之正確標準\ix{terms}{02@\zhlidxcat{正當性與合理歧見}!程序獨立之正確標準}（procedure-independent） & 根本法在表達並促進對每位公民平等尊重與關懷的道德要求時為正確 & 包含但不限於 Rawls 所稱 perfect 與 imperfect procedural justice，這兩種的正確性概念指向結果（例如兩人分得同樣大的蛋糕），不指向任何程序 \\
\bottomrule
\end{longtable}

註 19 說明，Michelman 的「pure-procedural」類別直接對應 Rawls 所稱「pure」程序正義；「procedure-independent」類別包含 Rawls 的「perfect」與「imperfect」程序正義，因這兩種的正確性概念都指涉結果，即使信奉該概念也帶來採用某程序的義務，如蛋糕例中「you cut, I pick」〔你切，我挑〕，原因在於該程序已知有較高的傾向達到所欲的結果。Michelman 在此處轉引 Rawls A Theory of Justice 85--86 頁。\footcite[405 n.19]{michelman1998}

訴諸程序獨立之政治正確標準，在 Post 看來與自治有深刻的衝突。Michelman 說這種直覺可以理解：根本法若須具有某些實質內容才正確，則這些法律在實質上應當如何，就不由任何人，更不由所有人自己決定；於是人民不能自治，至少不能正當地自治。\footcite[405--406]{michelman1998}

\section{Brennan 之判決與純粹程序民主觀}
\label{sec:20_responsive-4}

\subsection{與純粹程序民主觀相容之判決}

Michelman 問，Brennan 的言論自由判決，是否指向建立純粹程序的回應式民主\ix{terms}{01@\zhlidxcat{民主理論}!回應式民主}理論所要求的那種絕對不受限制之公共論述\ix{terms}{01@\zhlidxcat{民主理論}!不受限制之公共論述}領域。他引 Texas v. Johnson\ix{cases}{02@美國!01@聯邦最高法院!1989Texas v. Johnson（1989）@Texas v. Johnson（1989）}, 491 U.S. 397 (1989)，並指出 Brennan 加入 Hustler Magazine, Inc. v. Falwell\ix{cases}{02@美國!01@聯邦最高法院!1988Hustler Magazine, Inc. v. Falwell（1988）@Hustler Magazine, Inc. v. Falwell（1988）}, 485 U.S. 46 (1988) 的法院意見（Rehnquist 首席大法官執筆），Post 在 Constitutional Domains 119--178 頁把此案置於他對言論自由判例之詮釋的中心。\footcite[406 n.20]{michelman1998}這樣的讀法「surely not an incongruous reading」〔當然不是不協調的讀法〕。\footcite[406]{michelman1998}註 21 補充，有人對 Brennan 最著名的言論自由判決 New York Times Co. v. Sullivan\ix{cases}{02@美國!01@聯邦最高法院!1964New York Times Co. v. Sullivan（1964）@New York Times Co. v. Sullivan（1964）}, 376 U.S. 254, 269, 274--275 (1964) 中 Meiklejohn\ix{persons}{meiklejohn@Meiklejohn, Alexander} 式的修辭感到憂慮；Post 與 Dworkin\ix{persons}{dworkin@Dworkin, Ronald} 都在 Meiklejohn「it serves the course of wise decisions」〔有助於明智決定〕的理由中，看到限制公共論述\ix{terms}{01@\zhlidxcat{民主理論}!公共論述}的缺口（Dworkin Freedom's Law 199--205 頁；Post Constitutional Domains 268--289 頁；Michelman 107 Ethics 707--708、710 頁）。\footcite[406 n.21]{michelman1998}

\subsection{純粹程序民主觀無法說明之判決}

另一方面，Michelman 認為 Brennan 的工作中有大量內容，激進的程序主義民主理論無法解釋\ix{terms}{08@\zhlidxcat{憲法解釋與司法自制}!解釋}。其憲法法學的核心「apparently」〔看來〕站著一個相當厚的實質觀念：「the inestimable value of the ever-redeemable dignity of the individual」〔個人可以隨時獲得救贖之人性尊嚴，其價值不可估量〕（註 22 引 Michelman 自己的 Super Liberal，77 Va. L. Rev. 1261, 1268--1269 (1991)）。這個觀念在 Brennan 手中蘊含一組內容充實、強烈奠基性的基本人權，尊重這些權利是政治權力正當性的前提；最戲劇性的是它推動了 Brennan 的死刑判例（註 23 引 Super Liberal 同處與 Brennan 的 In Defense of Dissents，37 Hastings L.J. 427, 436 (1986)）。\footcite[406 nn.22--23]{michelman1998}Michelman 並列出由該觀念啟發的其他詮釋主張，各附判決於註 24 至 30：

\begin{longtable}{@{}L{0.36\linewidth}L{0.07\linewidth}L{0.49\linewidth}@{}}
\caption{Michelman 406--407 頁所列 Brennan 以人性尊嚴為核心之憲法主張及註 24--30 所引判決。}\\
\toprule
主張（Michelman 之概括） & 註 & 判決（卷頁） \\
\midrule
\endfirsthead
\toprule
主張（Michelman 之概括） & 註 & 判決（卷頁） \\
\midrule
\endhead
憲法保障請求政府給付之資格者受個別化聽證 & 24 & Goldberg v. Kelly\ix{cases}{02@美國!01@聯邦最高法院!1970Goldberg v. Kelly（1970）@Goldberg v. Kelly（1970）}, 397 U.S. 254 (1970) \\
擴張的不自證己罪與雙重危險之權利 & 25 & Malloy v. Hogan\ix{cases}{02@美國!01@聯邦最高法院!1964Malloy v. Hogan（1964）@Malloy v. Hogan（1964）}, 378 U.S. 1 (1964) \\
受協助而自願結束生命之權利 & 26 & Cruzan v. Director, Mo. Dep't of Health\ix{cases}{02@美國!01@聯邦最高法院!1990Cruzan v. Director, Missouri Department of Health（1990）@Cruzan v. Director, Missouri Department of Health（1990）}, 497 U.S. 261, 301--330 (1990)（Brennan 大法官不同意見） \\
性與生育自我決定之權利 & 27 & Eisenstadt v. Baird\ix{cases}{02@美國!01@聯邦最高法院!1972Eisenstadt v. Baird（1972）@Eisenstadt v. Baird（1972）}, 405 U.S. 438 (1972) \\
親密結社之自由 & 28 & Roberts v. United States Jaycees\ix{cases}{02@美國!01@聯邦最高法院!1984Roberts v. United States Jaycees（1984）@Roberts v. United States Jaycees（1984）}, 468 U.S. 609, 612--623 (1984) \\
在自己家中閱讀所欲閱讀之物 & 29 & Stanley v. Georgia\ix{cases}{02@美國!01@聯邦最高法院!1969Stanley v. Georgia（1969）@Stanley v. Georgia（1969）}, 394 U.S. 557 (1969)（Marshall\ix{persons}{marshall@Marshall, Thurgood} 大法官之意見） \\
使受刑人得反覆就定罪之實體與程序提出爭執 & 30 & Fay v. Noia\ix{cases}{02@美國!01@聯邦最高法院!1963Fay v. Noia（1963）@Fay v. Noia（1963）}, 372 U.S. 391 (1963) \\
\bottomrule
\end{longtable}

死刑一項在註 23 只引 Michelman 與 Brennan 的文章，未列判決。註 26 與註 29 標明撰寫意見者，其餘各註只列判決名稱與卷頁，Brennan 在各案中的角色，註文未交代。\footcite[406--407 \& nn.24--30]{michelman1998}

Michelman 的結論：這些立場無法由任何純粹程序的民主政府理想說明；純粹程序的理想也無法接受，法官在這類事項上的話應當壓過人民的反對。\footcite[407]{michelman1998}

\section{Post 理論所含之實質前提}
\label{sec:20_responsive-5}

Michelman 接著問，Post 的回應式民主\ix{terms}{01@\zhlidxcat{民主理論}!回應式民主}，是否如他所說的那樣純粹程序性。Michelman 的答覆是，它至少涉及某些實質：回應式民主理論把政治中個人自治的可能性建立在一項「loaded with normative content」〔承載規範內容〕的原則上，即每個人與每種觀點都絕對不受限制地接近公共論述\ix{terms}{01@\zhlidxcat{民主理論}!公共論述}。這項原則雖「about a procedural matter」〔關於程序事項〕，本身卻是實質事項；它揭示一項正當政府的基礎規範，持此說的人必須視之為不容辯論。\footcite[407]{michelman1998}

Michelman 指出，正因為理論以這項原則為基礎，Post 才得以用理論支持某些憲法詮釋的正確性勝過另一些，首先是對公共論述之內容或強度的最輕微限制也嚴格堅守的詮釋（註 31 指 Post Constitutional Domains 119--178 頁）。\footcite[407 n.31]{michelman1998}若回應式民主理論確有基礎預設，它便是 Brennan 憲法理論的「closer cousin」〔更近的表親〕，比人們原先所見更近；但這又回到原來的問題：任何以不容協商、不容辯論的規範預設為基礎的正當政府觀念，即使預設是像不受限制之公共論述\ix{terms}{01@\zhlidxcat{民主理論}!不受限制之公共論述}這種表面上屬於結構或形式的東西，是否能回答所有人如何真正在立法領域自治。\footcite[407]{michelman1998}

\section{仇恨言論之例與最高之法}
\label{sec:20_responsive-6}

Michelman 以 Post 的一項論證說明這個問題。美國的根本法包含對所有人保證「the freedom of speech」〔言論自由〕的條款，大多數美國人會贊同此條款為正確；爭議出在條款的適用，近例是公立大學校園或其他公共場所中對惡毒種族主義言論的憲法保護。\footcite[407 n.32]{michelman1998}可假設人人都忠誠地相信，政府限制種族上惡毒的言論，究竟是被禁止、被要求或被允許，有一個正確答案\ix{terms}{08@\zhlidxcat{憲法解釋與司法自制}!正確答案}；問題在於人們真誠地分歧，且在實際的政治時間內無法解決。\footcite[407--408]{michelman1998}

既然無法迴避此問題，就須有人違反反對者的意見而決定，在美國的制度中，這個人就是 Brennan 這樣的人。Michelman 問，這時「everyone's self-government」〔所有人之自治〕何在。除非 Brennan 的決定在充分說明之後，能被所有人看成直接源於某個更崇高、為所有人認可為正確的 metaconstitutional principle〔後設憲法原則\ix{terms}{03@\zhlidxcat{憲政主義與憲法之效力基礎}!後設憲法原則}〕，否則受 Brennan 之決定所控制的根本法，將「under guise of interpretation」〔以解釋為名〕由 Brennan 寫成，不由所有人寫成。\footcite[408]{michelman1998}Post 為此需要的後設憲法原則提出一項內容：保證所有人與所有觀點都能無偏見地接近不受限制之公共論述\ix{terms}{01@\zhlidxcat{民主理論}!不受限制之公共論述}，國族認同\ix{terms}{01@\zhlidxcat{民主理論}!國族認同}在其中不斷鍛造。Michelman 稱這項保證為「Really High Law」〔最高之法\ix{terms}{03@\zhlidxcat{憲政主義與憲法之效力基礎}!最高之法}〕，所有人皆可為它負制定之責。Brennan 以最高之法決定文本憲法這項「intermediately high law」〔中等高度的法〕，即言論自由條款，是否禁止、要求或允許政府限制種族言論。這使 Brennan 所決定的問題成為法之適用，不屬立法。\footcite[408]{michelman1998}

Michelman 緊接著問：「Or does it?」〔果真如此嗎？〕Brennan 不可能無爭議地適用，因為有些人主張種族言論本身限制了某些人把公共論述\ix{terms}{01@\zhlidxcat{民主理論}!公共論述}當作自治媒介的機會，不分享該主張所基於之生活經驗的人不易將之撥開（註 33 指 Post Constitutional Domains 319--320 頁）。\footcite[408 \& n.33]{michelman1998}Michelman 進而指出，這個關於最高之法如何適用的問題，與它的「identity」〔同一性〕，即其定義或內容，不能分開。「不受限制之公共論述」這個概念，在對言論自由條款適用於政府限制種族言論所作的不同詮釋之下，不能說仍是反映同樣道德、形上或意識形態承諾的同一個概念；而這項詮釋問題涉及道德與意識形態上負載的實踐問題，在今日的社會中爭論激烈。因此並非人人都同意同一部最高之法，至少不是在使人人藉為該部法之作者而自治的意義上同意；不受限制之公共論述原則作為最高之法，「either too indeterminate or too partisan」〔不是太不確定，就是太有黨派性〕，不足以構成所有人自治的實現。\footcite[409]{michelman1998}

\section{憲政民主之悖論與無窮後退}
\label{sec:20_responsive-7}

Michelman 把這個結果稱為「the paradox of constitutional democracy」〔憲政民主之悖論\ix{terms}{03@\zhlidxcat{憲政主義與憲法之效力基礎}!憲政民主之悖論}〕，並說它看來可以一般化。悖論的一面，是人們被迫把立法共識的可能性，亦即透過法律所有人自治的可能性，安置於「higher」〔較高〕的立法層次，因為根本法不涉及普通立法無法共識地解決的具體利益、信念與評價衝突。希望是每個人都能自主地判斷立法體系公平，因而把自己算作即使是自己深切反對的具體結果之制定者。回應式民主\ix{terms}{01@\zhlidxcat{民主理論}!回應式民主}的「所有人之自治」理論即以這種希望為前提，Post 的話是公民能「embrace the government as rightfully 'their own'」〔擁抱政府為正當地『屬於自己』〕，因為他們接受透過絕對不受限制之公共論述\ix{terms}{01@\zhlidxcat{民主理論}!不受限制之公共論述}來形成國族認同\ix{terms}{01@\zhlidxcat{民主理論}!國族認同}的體系是公平的，這個體系的維持灌輸公民「a sense of participation, legitimacy, and identification」〔參與、正當性與認同感〕（註 34 指 Constitutional Domains 273 頁）。\footcite[409 \& n.34]{michelman1998}

另一面是，根本法從具體爭議中抽象出來，製造了一個問題：解釋這些抽象法律所要解決的事項，通常本身就是重大的政治道德問題，這些問題的解決方式與根本法本身的決定無法分開。為此，追求個人政治自治的要求，不僅初步的抽象表述須如此完成，所有重大的詮釋也須如此；而取得有效力的、類似法律的權威詮釋，只能經由某種非全體一致的集體過程，以制度方式進行。唯一能使這些詮釋與所有人之自治相容的方式，是在受詮釋的根本法之後，另有一組更抽象、更根本的原則，即最高之法\ix{terms}{03@\zhlidxcat{憲政主義與憲法之效力基礎}!最高之法}，可歸屬於所有人之制定，同時能客觀地決定相互競爭的根本法詮釋。但最高之法比根本法更抽象，在決定所需事項的過程中就更確定地需要有爭議的詮釋。Michelman 稱這是一個「infinite regress」〔無窮後退\ix{terms}{03@\zhlidxcat{憲政主義與憲法之效力基礎}!無窮後退}〕。\footcite[409--410]{michelman1998}註 35 把它與批判法學者對所謂「liberal legalism」的反對相比，列 Michelman 自己的 Justification (and Justifiability) of Law in a Contradictory World，載於 Pennock 與 Chapman 編 Justification 71, 79--81 頁，該文所述為 Duncan Kennedy 於 28 Buffalo L. Rev. 205 (1979) 的論證。\footcite[410 n.35]{michelman1998}

Michelman 說這可能只是以冗長的方式反對一切「non-foundational」〔非奠基式〕的憲法正確性觀念。憲政主義\ix{terms}{03@\zhlidxcat{憲政主義與憲法之效力基礎}!憲政主義}，即「the endeavor to place government under reason expressed as law」〔使政府服從於以法律表達之理性的努力〕，似乎必然意味著建立某組先驗固定、不容協商、不容辯論的規範前提。就目前的目的而言，把它們稱為權力正當性的前提，或稱為決定權力正當性之辯論的正當性之前提，並無差別。例子有絕對不受限制之公共論述\ix{terms}{01@\zhlidxcat{民主理論}!公共論述}，以及對個人隨時可獲救贖之人性尊嚴的堅定尊重。這些可視為憲政民主的文化承諾，人們因而被推向一個結論：只有在一國人民普遍先驗地接受它們這種幸運的情況下，自治與憲政才可能並存。\footcite[410]{michelman1998}Michelman 說，在這點上 Brennan 的生涯指向正確、也不可避免的方向；但在人民對正當政府基本規範原則有先驗一致的這項前提之下，決定根本法時是否還有民主的位置，就不再明顯。註 36 以 Rawls 較晚近的政治理論相比：該理論賦予重要地位的，是某些國家公共政治文化中潛伏的一種觀念，即把每個人視為政治上自由而平等者（Mulhall 與 Swift 的 Liberals and Communitarians 189--190 頁；Rawls 的 Political Liberalism 13--14 頁）。\footcite[410 \& n.36]{michelman1998}

\section{與 Post 438 頁註 33 之互見}
\label{sec:20_responsive-8}

Michelman 所稱 Post 的純粹程序主張，Post 在 438 頁註 33 有所回應：他說 Michelman 歸給他的程序性民主說明，大致只表達了民主之實質觀念對公共論述\ix{terms}{01@\zhlidxcat{民主理論}!公共論述}之規範的一項蘊含（註引 Constitutional Domains 277--278 頁）。\footcite[438 \& n.33]{post1998}Post 在同頁正文把人民主權\ix{terms}{01@\zhlidxcat{民主理論}!人民主權}界定為依程序標準，把民主界定為依實質價值，相關討論見第 \ref{ch:70_misreadings} 部。
